\chapter{Literature Review}
\label{ch:2}

\section{Preliminaries}
\label{sec:2.1}

\textbf{Empathy as something that can be measured:} a system cannot be built against a vague word. This work uses EPITOME for text-based mental-health support, built on well-known empathy scales, which divides a reply into \textit{emotional reactions}, naming and validating the feeling; \textit{interpretations}, showing the situation has been understood; and \textit{explorations}, inviting the person to say more \cite{sharma2020epitome}. Counselling practice adds the opposite limit, warning of the question-and-answer trap in which the person becomes a passive respondent \cite{samhsa2019tip35}. Together they give a working definition a machine can be held to: acknowledge the feeling, respond to the situation, and invite without interrogating.

\textbf{The corpus:} EmpatheticDialogues is the standard benchmark for this task \cite{rashkin2019}. One crowdworker was given one of 32 emotion labels, wrote a description of a matching situation and then described it in conversation, while a second replied as a listener without seeing the label. The released files hold 24,850 conversations split so that no speaker appears in more than one set \cite{ed_github}. Two features matter here: the situation description and the first message are distinct texts, the description written first; and each conversation carries exactly one label, chosen before it was composed.

\textbf{GoEmotions }is the other dataset commonly used for fine-grained emotion: about 58,000 Reddit comments labelled with 27 emotions plus neutral, allowing more than one label per comment \cite{demszky2020}. This project began there and moved to EmpatheticDialogues for four reasons measured in Section 4.3.1.

\textbf{Emotion recognition in conversation} is a field of its own rather than sentence-level classification, because the same words carry different feelings depending on what preceded them \cite{poria2019erc}. Related work classifies the intents a listener's reply can express on the same corpus \cite{welivita2020}, and shows that emotions in dialogue are often implied and must be inferred rather than read off \cite{koga2023implicit}, and that one reply can carry more than one label \cite{belay2025}.

\textbf{Models:} every model here derives from the transformer \cite{vaswani2017}, in two families doing two jobs. \textbf{Encoders }are trained with masked-language modelling, which suits them to classification; BERT set the precedent \cite{devlin2019} and RoBERTa extended it with more text, longer training and dynamic masking \cite{liu2019roberta}. \textbf{Decoder-only language models} generate text left to right and, when instruction-tuned, follow written instructions; we use Qwen2.5 \cite{qwen2024} and Phi-4-mini \cite{abouelenin2025phi4mini}, the latter trained with a particular emphasis on instruction following.

\textbf{Two techniques bring large models to size} on a consumer GPU. \textbf{Quantisation }stores weights at reduced precision; 4-bit NF4 with double quantisation keeps the quality of 16-bit weights at roughly a quarter of the memory \cite{dettmers2023qlora}. \textbf{Low-rank adaptation} trains a small set of extra parameters rather than the whole network, making fine-tuning feasible on the same hardware \cite{hu2022lora}. A third, \textbf{retrieval-augmented generation}, supplies relevant passages retrieved at inference time \cite{lewis2020rag}, usually through sentence embeddings \cite{reimers2019sbert}; this work uses it as a baseline but not in the final design.

\textbf{Measure of empathy:} no single automatic measure captures empathy, so three kinds of evaluation are used here. The\textbf{ WASSA shared task on empathy detection }provides a model trained on human empathy ratings of conversation turns \cite{giorgi2024wassa}, used here as the hyy33 system \cite{yang2024hyy33}.\textbf{ Large language models as judges} track human preference well when the judge is strong \cite{zheng2023judge} but favour their own generations \cite{panickssery2024}, so a panel from different developers is recommended over a single judge \cite{verga2024}. The reference remains \textbf{human evaluation}, reported either as a five-point rating \cite{qian2023llm} or a head-to-head preference \cite{cao2025tooled}.

\section{Review of Existing Research}
\label{sec:2.2}

The relevant work falls into six groups: the benchmark that began the field; the models trained end to end on it; work on the listener rather than the speaker; the arrival of large language models, which changed what counts as a good reply; the narrower but decisive question of whether those models recognise emotion or only write about it well; and the clinical work using the same machinery for screening rather than replying. Two further sections cover the study closest to ours and how any of this should be judged.

\subsection{The Benchmark That Created the Field}
\label{sec:2.2.1}

Rashkin et al. (2019) created the field by releasing EmpatheticDialogues with a framing every later paper has followed: a listener answers a person's description of an emotional situation, and the emotion is one of 32 fine-grained categories rather than a sentiment \cite{rashkin2019}. That choice is why top-1 accuracy here sits far below figures on two- or six-class tasks: it forces a system to separate feelings a sentiment model would merge, such as \textit{sad}, \textit{lonely}, \textit{disappointed }and \textit{devastated}.

Two further decisions in that paper shaped everything after it. The emotion label may either be given to the generator or predicted by it, and the paper shows both. And it set a reference point for recognition: a fastText classifier at about 43\% accuracy \cite{rashkin2019}. That figure is widely cited, and it was measured on the situation description, the text the crowdworker wrote about the event, not on the first message the speaker sends the listener. The two differ in length and explicitness, since a description of an event names the feeling more openly than a message written to another person. Any system working on what a user actually types is solving the harder of the two problems, and Section 5.2.2 gives our own measurement of that gap.

\subsection{Joint Models Trained End to End}
\label{sec:2.2.2}

A series of \textbf{joint models} followed, each adding a mechanism to a single encoder-decoder network trained from scratch on the corpus. MoEL (Lin et al., 2019) used a mixture of listeners, one per emotion, fused by a softmax over the detected emotion so the model could answer different feelings differently \cite{lin2019moel}. MIME (Majumder et al., 2020) started from an observation about human listeners, that they partly mirror the speaker's emotion and partly compensate for it \cite{majumder2020mime}. EmpDG (Li et al., 2020) added an adversarial discriminator judging whether a reply looked empathetic, together with emotion signals at more than one granularity \cite{li2020empdg}. The next three brought in knowledge the corpus does not contain. KEMP (Li et al., 2022) added external commonsense and concept knowledge, letting the model infer about the situation rather than only its words \cite{li2022kemp}. CEM (Sabour et al., 2022) directed those inferences at the speaker, eliciting their likely reactions, intentions and needs \cite{sabour2022cem}. CASE (Zhou et al., 2023) aligned a commonsense graph with a concept graph, on the view that a listener must relate what the speaker thinks to what they feel \cite{zhou2023case}. ESCM (Yang et al., 2023) then used emotion-semantic vectors that change across the conversation, supervised explicitly on the emotional signal \cite{yang2023escm}.

Read in order, these papers show emotion accuracy creeping up: 30.6\% for EmpDG, 30.8\% for MoEL, 36.6\% for KEMP, 37.8\% for CEM, 38.9\% for CASE and 41.2\% for ESCM. Two things follow. The gains come from added knowledge and structure rather than from what is done with the emotion once predicted; and after four years of architectural work the best of them still misreads the speaker's feeling in about three cases out of five. A related model suggests why. Liu et al. (2021) ran RoBERTa-GPT2 on the same task and corpus but initialised its encoder from pre-trained RoBERTa and its decoder from pre-trained GPT-2 instead of training from scratch, reaching 51.5\% \cite{liu2021robertagpt2}, about ten points above the best knowledge-augmented model here. On a corpus this size the limiting factor was not what was built into the models but how little text they had seen, which is the first reason our design spends its parameter budget on strong pre-trained components and on how they are connected.

\subsection{The Listener's Side of the Conversation}
\label{sec:2.2.3}

Most of the above models the speaker. A smaller body of work looks at the listener instead, and for a rule-based design it is more useful than its citation count suggests. From a large sample of EmpatheticDialogues, Welivita and Pu (2020) built a taxonomy of empathetic response intents showing that the listeners in this corpus use a small, describable set of behaviours: questioning, acknowledging, consoling, agreeing, encouraging, sympathising and a few more \cite{welivita2020}. That matters for any rule-writing approach, ours included, because it shows supportive replies have structure. If human listeners work from a short list of moves, asking a generator to name the feeling, answer the situation and ask no more than one question is not an arbitrary constraint but a description of what the corpus already does.

The wider literature on emotion recognition in conversation supplies the other half. Poria et al. (2019) explain why it is not sentence-level classification: the same words carry different feelings depending on what precedes them, so context is an input to the system rather than an optional extra \cite{poria2019erc}. That is why how much conversation the detector reads is treated here as a parameter to be measured rather than assumed, as reported in Section 5.3.2.

\subsection{Large Language Models Enter the Task}
\label{sec:2.2.4}

\textbf{Large language models} changed the field, and Qian et al. (2023) made that plain. Their human evaluation of ChatGPT against the purpose-built models above found raters overwhelmingly favouring ChatGPT, 4.64 out of 5 against 3.08 for CEM \cite{qian2023llm}, a gap tuning cannot explain. The paper then took a path directly relevant to our design: replies improve when the emotion is given explicitly, and improve further when the model is told to reason about the situation before answering. The best system in this literature was therefore a general model given structure from outside, which is what our two-stage pipeline supplies. Loh and Raamkumar (2023) reached a similar conclusion against counselling criteria for online mental-health support, adding that without structure the replies stay generic \cite{loh2023harnessing}.

That structure takes several forms, each with a price. TOOL-ED (Cao et al., 2025) treats a commonsense knowledge base as a tool the model may call during generation, deciding when to consult COMET or CICERO rather than having their output injected into every input, which avoids the noise unrestricted knowledge introduces and frees the system from any one knowledge base \cite{cao2025tooled}. It is the knowledge built into KEMP and CEM, now pulled as needed. The cost shows in the setup: a fine-tuned 7B model trained on 48 GB GPUs \cite{cao2025tooled}, six times our GPU memory.

Wu et al. (2025) argue the same situation should be answered differently for different people, conditioning generation on an MBTI classifier of about 84\% accuracy and on images alongside the text \cite{wu2025traits}. Their model reaches 45.4\% emotion accuracy, dropping almost 4 points without the personality signal \cite{wu2025traits}. The constraint for our aim is the input: the design assumes visual data and personality labels, and a text-only system reading a single message has neither.

Wang et al. (2025) entered the NLPCC 2025 emotional support conversation task, where the target is a strategy sustained across a conversation, comparing low-rank adaptation with full fine-tuning on Qwen2.5 \cite{wang2025esc}. Two facts apply here. Their best configuration raised the combined score from 30.87 to 41.17 using low-rank adaptation, with the adapter's rank and scaling factor affecting the score \cite{wang2025esc}, and their settings bracket ours (Section 4.4.4). The other is scale: the entry fine-tuned models up to 72B parameters on 1,500 dialogues, the opposite end of the resource question from the system built here.

Rasool et al. (2025) work on psychotherapy transcripts rather than crowdsourced conversations, injecting emotion-aware embeddings at word, sentence and session level into Flan-T5, Llama 2 13B, DeepSeek-R1 and GPT-4, with a vector index retrieving related passages over those embeddings \cite{rasool2025}. The retrieval part resembles the baseline this thesis builds from sentence embeddings at generation time \cite{reimers2019sbert,lewis2020rag}, and is a useful comparison for Section 5.2.5, where our retrieval variant is measured and rejected. They report gains in empathy, coherence and fluency over their baselines, but note that their paired significance tests fall short of the conventional threshold and that the study has no human evaluation \cite{rasool2025}. Li et al. (2024) take the emotional route instead of the knowledge one, having the model reason step by step about emotion before answering \cite{li2024emotioncot}; our stage-one emotion choice is a condensed form of the same idea, one deliberate choice, recorded, before any reply is written.

\subsection{Do These Models Recognise Emotion?}
\label{sec:2.2.5}

All of that concerns the quality of the answers. Whether large models reliably \textit{detect }emotion is a separate question, and the one this thesis turns on, since a decoupled design is only worth building if a dedicated detector is genuinely better. The evidence says it is.

Feng et al. (2023) tested large language models on affect recognition in conversation and found them comparable to, not better than, purpose-trained classifiers \cite{feng2023affect}. Belay et al. (2025) settled the comparison, building a multi-label emotion dataset in four Ethiopian languages (EthioEmo), adding the English SemEval-2018 data, and comparing encoder-only, encoder-decoder and decoder-only models together with in-context learning. Their conclusion is plain: fine-tuned smaller pre-trained models beat the zero-shot performance of large language models on this task, the best decoder-only averaging about 46 against about 72 for a fine-tuned encoder on its best language \cite{belay2025}. They also find that correct multi-label classification remains hard even in English, and that well-resourced and poorly resourced languages differ sharply \cite{belay2025}, a point we return to in Section 5.6.5, since our system is English only.

Two further papers explain \textit{why }recognition is hard rather than \textit{how well} it is done. Koga et al. (2023) studied emotions that are evoked but not expressed and concluded that much of the problem is inference rather than classification \cite{koga2023implicit}. The second difficulty is Belay et al.'s own, that an utterance often carries more than one emotion at once \cite{belay2025}, which is one reason our detector returns three candidates instead of one.

The best single figure comes not from a benchmark paper but from a deployed system. The Chirpy Cardinal agent built for the Alexa Prize carried a RoBERTa classifier over the same 32 categories, reported at about 61.5\% accuracy on the first utterance, with errors concentrated on neighbouring emotions rather than scattered \cite{paranjape2020chirpy}. Our results reproduce both parts: Section 5.1.2 reports the accuracy and the same neighbour-heavy error pattern, which is what makes a three-candidate output useful rather than evasive, and Section 5.5.1 explains the difference between the two figures.

\subsection{Clinical and Mental-Health Applications}
\label{sec:2.2.6}

The clinical end of the field uses the same machinery for screening rather than replying, and sets the standard of care such a system will be held to. Lan et al. (2024) applied large language models to depression detection on social media \cite{lan2024depression}, and the CLPsych 2025 shared task recast the problem as tracking mental-health dynamics across a user's timeline rather than classifying a single post \cite{tseriotou2025clpsych}. Kumar et al. (2024) survey the two problems limiting the area: labelled clinical data is scarce, and the bias from training on what little exists is hard to detect \cite{kumar2024unlocking}.

The caution most relevant to our design comes from the systematic review by Sorin et al. (2024), which we act on directly. That large language models were often scored as empathetic as clinicians or more so, across the studies it examined, is remarkable in itself. The problem the review identifies is not incompetence but sameness: the empathy was repetitive, the same formulations recurring across answers \cite{sorin2024jmir}. A comparison with a reference cannot detect that failure, since a repetitive reply may sit very close to a good one. Hence the repetition measure among our own signals in Section 5.1.4, which none of the joint models above reports.

\subsection{The Closest Prior Work}
\label{sec:2.2.7}

The study most similar to this thesis is the \textbf{Hybrid Empathetic Framework} by Yang et al. (2024). It first adds a small, specialised empathetic model to a large language model, and then predicts the emotion in two steps, along with determining the reason for the emotion \cite{yang2024hef}. That's the same architectural intuition that this thesis seeks to achieve: a small model that is good at emotion feeding a large model that is good at language, with the emotional decision made in a stage of its own, not folded into generation.

The systems differ in purpose. Their framework targets benchmark performance on empathetic generation. Ours runs entirely within 8 GB on a consumer GPU so a user's disclosure never leaves the machine; exposes the whole emotional decision, the candidates, their confidence and the words behind them, as an audit trail on every reply; and judges replies by what a supportive reply must do rather than by how closely they match a reference. Because that paper exists, our contribution is not the idea of separating detection from generation, but what the separation provides and the evidence gathered for it.

Table~\ref{tab:2.1} compares the works that are most similar to this thesis side by side, what each used, what each accomplished, and the limitation that is pertinent here.

\begin{table}[htbp]
\centering
\footnotesize
\caption{Works that are most similar to this thesis and limitation of each work}
\label{tab:2.1}
\begin{tabularx}{\textwidth}{@{}lXXX@{}}
\toprule
\textbf{Work} & \textbf{Data} & \textbf{Approach and result} & \textbf{Limitation relevant to this thesis} \\
\midrule
Rashkin et al., 2019 (ACL) & EmpatheticDialogues & Introduces 25k-conversation, 32-emotion benchmark; emotion classifiers on situation text & Benchmark only; situation-text classifiers overstate real-input accuracy \\
Lin et al., 2019 - MoEL (EMNLP) & ED & Mixture of emotion-specific listeners; emotion acc 30.75\% & Low recognition accuracy; small model generation \\
Majumder et al., 2020 - MIME (EMNLP) & ED & Emotion mimicry grouped by polarity & Coarse polarity groups; generic replies \\
Sabour et al., 2022 - CEM (AAAI) & ED & Adds COMET commonsense (cognitive + affective) & Human empathy 3.08/5, far below LLMs \\
Zhou et al., 2023 - CASE (ACL) & ED & Aligns cognition and affection at coarse and fine grain & Human empathy 3.04/5 \\
Qian et al., 2023 (arXiv 2310.05140) & ED & LLMs (ChatGPT) beat trained SOTA on automatic and human metrics & Uncontrolled, verbose; no rule/format control measured \\
Paranjape et al., 2020 - Chirpy Cardinal & ED + live Alexa users & RoBERTa-large emotion classifier on first utterance, 61.5\% top-1 & Technical report; validation-set score \\
Cao et al., 2025 - TOOL-ED (COLING) & ED & LLM tool-calling for commonsense knowledge; A/B human eval & Fine-tuning 7-14B models; heavy resources \\
Wu et al., 2025 - Personality-aware ERG (COLING) & MELD-based multimodal + ED & Personality + visual cues; emotion acc 45.4\% text-only & Needs multimodal input and personality labels \\
Belay et al., 2025 (COLING) & EthioEmo (4 Ethiopian languages) + SemEval-2018 English & Encoder, encoder-decoder, decoder LLMs for multi-label emotion & Fine-tuned encoders still beat few-shot LLMs at emotion classification \\
Rasool et al., 2025 (AI, MDPI) & Psychotherapy transcripts & Emotion-aware embedding fusion into Flan-T5/Llama 2/DeepSeek-R1/GPT-4; BERT for emotion & Paired t-tests non-significant; no human evaluation \\
Wang et al., 2025 (arXiv 2507.12820) & NLPCC 2025 ESC task & Prompting + LoRA / full fine-tuning of Qwen2.5; total 41.17 vs 30.87 base, ranked 2nd & Needs fine-tuning up to 72B parameters \\
\bottomrule
\end{tabularx}
\end{table}

\subsection{How Empathetic Replies Should Be Judged}
\label{sec:2.2.8}

Finally, \textbf{how this should be measured} has become a literature of its own, and it is where this thesis departs furthest from the earlier papers, since our question can only be answered by measurement.

Zheng et al. (2023) found strong language models used as judges agreeing with human preference over 80\% of the time, comparable to agreement between people, and documented their shared biases, such as a preference for longer answers \cite{zheng2023judge}. Panickssery et al. (2024) showed that such evaluators recognise and favour their own generations, making the choice of evaluator a methodological decision rather than a convenience, since a judge from the same family as a system under test is not a neutral instrument \cite{panickssery2024}. The remedy this thesis adopts, a panel of judges from different model families, was proposed by Verga et al. (2024) on grounds of both reduced bias and lower cost \cite{verga2024}.

The complementary solution is to train a model directly on human empathy ratings rather than have it imitate a rater \cite{giorgi2024wassa}, as in the hyy33 system used here \cite{yang2024hyy33}; its value is being unrelated to every system evaluated. Human evaluation remains the benchmark, reported recently either as a five-point rating \cite{qian2023llm} or a head-to-head comparison \cite{cao2025tooled}. This thesis uses all three, with the judges and the classifier in one panel, standardised before averaging (Section 5.1.5).

\section{Summary of Key Findings}
\label{sec:2.3}

The review yields seven findings, each stated with its evidence and with what this thesis does about it.

\textbf{Pre-training matters more than architecture for this task:} four years of joint models moved emotion accuracy from 30.6\% to 41.2\% \cite{lin2019moel,majumder2020mime,li2020empdg,li2022kemp,sabour2022cem,zhou2023case,yang2023escm}, while RoBERTa-GPT2 reaches 51.5\% on the same task simply by using pre-trained components instead of training from scratch \cite{liu2021robertagpt2}. Ten points separate them and the cheaper method wins, so the parameter budget is better spent on a strong pre-trained model, and on how it is connected, than on a network built from 25,000 conversations.

\textbf{Large language models write better replies than purpose-built empathy models but do not reliably detect emotion: }human raters prefer ChatGPT to CEM, 4.64 against 3.08 \cite{qian2023llm}, and a systematic review finds large models rated level with or above clinicians on empathy \cite{sorin2024jmir}. Yet they do no better than trained classifiers on conversational affect \cite{feng2023affect}, worse than fine-tuned encoders on multi-label emotion \cite{belay2025}, and struggle with implied emotion \cite{koga2023implicit}. That asymmetry is the direct argument for a separate detector and a generator doing only what it is good at.

\textbf{A small fine-tuned model can beat a large prompted one at classification: }the sharpest evidence for our architecture comes from outside the empathy literature. Belay et al. found fine-tuned smaller models still beating the zero-shot performance of large language models on the same multi-label dataset \cite{belay2025}, and our own corpus has a matching figure, a RoBERTa classifier at about 61.5\% on the 32 EmpatheticDialogues emotions \cite{paranjape2020chirpy}. A detector of a few hundred million parameters is therefore not a hardware compromise but, on this task, the better instrument.

\textbf{Machine empathy is repetitive and hardly anyone measures it: }the weakness the systematic review identifies is not incompetence but sameness, fluent and kind replies drawn from a small pool of phrases \cite{sorin2024jmir}. None of the joint models above reports a repetition measure, and a reference-based comparison cannot find the problem, since a repeated reply may sit close to the reference. We therefore report repetition for every system, not only ours.

\textbf{Structure helps, and there are several ways to supply it, each with a cost: }telling the model the emotion or asking it to reason before answering \cite{qian2023llm,li2024emotioncot}; offering a knowledge base as a callable tool \cite{cao2025tooled}; conditioning on inferred personality and visual input \cite{wu2025traits}; fine-tuning with a low-rank adapter \cite{wang2025esc}; or retrieving related passages at generation time \cite{lewis2020rag,rasool2025}. All report gains, but none establishes what is worth its cost under a hardware constraint: tool calling needs 48 GB \cite{cao2025tooled}, the fine-tuned approach models up to 72B parameters \cite{wang2025esc}, and the trait-aware approach multimodal input with personality labels \cite{wu2025traits}. Chapter 5 answers that within 8 GB, including the retrieval variant, which we measure and dismiss.

\textbf{Evaluation is the field's weakest aspect, and workable alternatives now exist: }published comparisons rest on overlap with a single reference reply, which cannot tell whether a reply acknowledged the feeling or asked five questions in a row, and one recent study reports improvements that miss the conventional threshold and that no human rated \cite{rasool2025}. The alternatives include a classifier trained directly on human empathy ratings \cite{giorgi2024wassa,yang2024hyy33} and panels of judges from different developers, which offset the self-preference any single judge shows \cite{zheng2023judge,verga2024,panickssery2024}. They let this thesis judge replies by what they \textit{do} rather than what they resemble.

\textbf{Separating a small model from a large one has been proposed, but for another purpose: }the Hybrid Empathetic Framework pairs a small empathetic model with a large one and predicts the emotion in two steps \cite{yang2024hef}. Our work shares that structure and differs in purpose: a whole system within 8 GB, the emotion exposed as an auditable record rather than an internal step, and an evaluation built on rule adherence, an empathy panel and paired statistics.

\textbf{The gap this thesis fills:} to our knowledge no existing work builds an empathetic response system that (i) runs entirely on consumer hardware, so nothing a user discloses leaves the machine, (ii) exposes its whole emotional decision for inspection before a reply is written, and (iii) is evaluated against single-model baselines of both its own size and a larger one, on measures matching what a supportive reply must do, with every difference tested statistically. Chapters 4 and 5 build and test that system.
